%% =================================================================
%% Beber et al. IEEE Robotics and Automation Letters, Feb 2026.
%% Reconstruction of page 8 (printed p. 8).
%% Contents: Fig. 5 placeholder (estimated stiffness distribution on
%% silicone phantom, 50 mm square domain, experimental trajectory);
%% Section VI Conclusion; References [1]--[29].
%% =================================================================

\documentclass[11pt]{article}

\usepackage[margin=0.85in]{geometry}
\usepackage{amsmath, amssymb}
\usepackage{mathrsfs}
\usepackage{xcolor}
\usepackage{fancyhdr}
\usepackage{url}

\pagestyle{fancy}
\fancyhf{}
\lhead{\small 8\ \ IEEE RAL, Accepted Feb 2026}
\rhead{\small }
\cfoot{\small 8 $\to$ \arabic{page}}
\renewcommand{\headrulewidth}{0.4pt}

\setcounter{page}{1}
\setcounter{equation}{5}
\setlength{\parskip}{6pt}
\setlength{\parindent}{0pt}

\begin{document}

% ---------- Fig. 5 placeholder ----------
\begin{center}
\fbox{\begin{minipage}[c][6cm][c]{0.90\textwidth}
\centering
\textbf{Fig. 5.}\ Estimated stiffness distribution from palpation of
the soft phantom.\\[4pt]
\textcolor{gray}{\footnotesize
50 mm $\times$ 50 mm domain in the $(x, y)$ plane. Colour scale:
estimated elastic modulus ranging from $\approx 102$ kPa (soft
background) up to $\approx 183$ kPa (stiff DragonSkin-30NV inclusion
shown in yellow). The end-effector trajectory during the search
motion is overlaid as a grayscale-gradient path: black marks the
beginning, white the end. The trajectory concentrates exploration in
the stiffer region on the left while still covering the rest of the
map, illustrating the adaptive focus of the EID-driven ergodic
planner.}
\end{minipage}}
\end{center}

\vspace{0.4em}

\textbf{Fig. 5.}\ Estimated stiffness distribution obtained through
palpation of the soft phantom. The stiffer region is shown in yellow.
The end-effector trajectory during the search motion is visualised
using a grayscale gradient, from black (beginning) to white (end). A
video of the experiment is provided in the Supplemental Materials
and is also available online at
\url{https://youtu.be/a3BUSxA7wXY}.

\vspace{0.8em}

% ---------- Section VI ----------
\section*{VI.\ Conclusion}

This work presents an autonomous robotic palpation framework for
real-time elastic stiffness mapping of soft tissues. By combining
force-based viscoelastic parameter estimation with ergodic trajectory
planning implemented via a heat-equation-driven controller, the
method enables continuous palpation without embedded tactile sensing.
Exploration is guided by an EID tailored to stiffness mapping,
allowing adaptive focus on diagnostically relevant regions.
Simulation results demonstrate improved stiffness reconstruction
accuracy, with reliable detection of multiple stiff regions and
lower RMSE and segmentation errors compared to BO, particularly
under noisy conditions. Experimental validation with a robotic arm
and a silicone phantom confirms the robustness of the approach,
showing consistent detection of stiff inclusions and good agreement
with ground-truth stiffness maps. Overall, the proposed ergodic
strategy effectively balances exploration and accurate estimation,
making it well suited for autonomous palpation as an assistive
diagnostic tool. Future work will extend the method to
three-dimensional mapping on non-flat surfaces and evaluate it on
anatomically realistic phantoms and ex vivo tissues.

% ---------- References ----------
\section*{References}

\begin{footnotesize}
\noindent
[1] R.~Tozzi, C.~K\"ohler, et al., ``Laparoscopic treatment of early
ovarian cancer: surgical and survival outcomes,'' Gynecologic
oncology, vol.~93, no.~1, pp.~199--203, 2004.\\[2pt]
[2] N.~Eisemann, S.~Bunk, et al., ``Nationwide real-world
implementation of AI for cancer detection in population-based
mammography screening,'' Nature Medicine 2025 31:3, vol.~31, no.~3,
pp.~917--924, 1 2025.\\[2pt]
[3] Q.~Dan, Z.~Xu, et al., ``Diagnostic performance of deep learning
in ultrasound diagnosis of breast cancer: a systematic review,'' npj
Precision Oncology 2024 8:1, vol.~8, no.~1, pp.~1--13, 1 2024.\\[2pt]
[4] A.~Attanasio, B.~Scaglioni, et al., ``Autonomy in Surgical
Robotics,'' Annual Review of Control, Robotics, and Autonomous
Systems, vol.~4, 2021.\\[2pt]
[5] L.~Zhang, Y.~Peng, et al., ``Innovative integration of automated
breast volume scan and ultrasound elastography for enhanced
differentiation of benign and malignant breast lesions,'' Scientific
Reports 2025 15:1, vol.~15, no.~1, pp.~1--8, 5 2025.\\[2pt]
[6] I.~Sack, ``Magnetic resonance elastography from fundamental
soft-tissue mechanics to diagnostic imaging,'' Nature Reviews
Physics, vol.~5, no.~1, 2023.\\[2pt]
[7] N.~Pacheco, C.~Henry, et al., ``Virtual palpation: A novel method
to identify the physical properties of tissue for laser surgery,''
in Hamlyn Symposium on Medical Robotics, 2025.\\[2pt]
[8] A.L.~Trejos, J.~Jayender, et al., ``Robot-assisted tactile
sensing for minimally invasive tumor localization,'' International
Journal of Robotics Research, vol.~28, no.~9, 2009.\\[2pt]
[9] N.~Enayati, E.~De~Momi, and G.~Ferrigno, ``Haptics in
robot-assisted surgery: Challenges and benefits,'' IEEE Reviews in
Biomedical Engineering, vol.~9, 2016.\\[2pt]
[10] S.~McKinley, A.~Garg, et al., ``An interchangeable surgical
instrument system with application to supervised automation of
multilateral tumor resection,'' in IEEE International Conference on
Automation Science and Engineering, 2016.\\[2pt]
[11] N.~Zevallos, R.~Arun~Srivatsan, et al., ``A Real-time Augmented
Reality Surgical System for Overlaying Stiffness Information,'' in
Robotics: Science and Systems, 2018.\\[2pt]
[12] A.~Garg, S.~Sen, et al., ``Tumor localization using automated
palpation with Gaussian Process Adaptive Sampling,'' in IEEE
International Conference on Automation Science and Engineering,
2016.\\[2pt]
[13] E.~Ayvali, R.A.~Srivatsan, et al., ``Using Bayesian optimization
to guide probing of a flexible environment for simultaneous
registration and stiffness mapping,'' in Proceedings - IEEE
International Conference on Robotics and Automation, 2016.\\[2pt]
[14] H.~Salman, E.~Ayvali, et al., ``Trajectory-Optimized Sensing for
Active Search of Tissue Abnormalities in Robotic Surgery,'' in
Proceedings - IEEE International Conference on Robotics and
Automation, 2018.\\[2pt]
[15] Y.~Yan and J.~Pan, ``Fast Localization and Segmentation of
Tissue Abnormalities by Autonomous Robotic Palpation,'' IEEE
Robotics and Automation Letters, vol.~6, no.~2, pp.~1707--1714,
4 2021.\\[2pt]
[16] E.~Ayvali, A.~Ansari, et al., ``Utility-Guided Palpation for
Locating Tissue Abnormalities,'' IEEE Robotics and Automation
Letters, vol.~2, no.~2, 2017.\\[2pt]
[17] K.A.~Nichols and A.M.~Okamura, ``Methods to Segment Hard
Inclusions in Soft Tissue During Autonomous Robotic Palpation,''
IEEE Transactions on Robotics, vol.~31, no.~2, 2015.\\[2pt]
[18] L.~Beber, E.~Lamon, et al., ``Towards Robotised Palpation for
Cancer Detection through Online Tissue Viscoelastic Characterisation
with a Collaborative Robotic Arm,'' in IEEE International Conference
on Intelligent Robots and Systems, 2024, pp.~2380--2386.\\[2pt]
[19] S.~Ivic, B.~Crnkovic, and I.~Mezic, ``Ergodicity-Based
Cooperative Multiagent Area Coverage via a Potential Field,'' IEEE
Transactions on Cybernetics, vol.~47, no.~8, 2017.\\[2pt]
[20] C.E.~Rasmussen and C.K.~Williams, Gaussian Processes for
Machine Learning. The MIT Press, 2006.\\[2pt]
[21] P.~Chalasani, L.~Wang, et al., ``Concurrent nonparametric
estimation of organ geometry and tissue stiffness using continuous
adaptive palpation,'' in Proceedings - IEEE International Conference
on Robotics and Automation, vol.~2016-June, 2016.\\[2pt]
[22] ---, ``Preliminary evaluation of an online estimation method
for organ geometry and tissue stiffness,'' IEEE Robotics and
Automation Letters, vol.~3, no.~3, 2018.\\[2pt]
[23] G.~Mathew and I.~Mezi, ``Metrics for ergodicity and design of
ergodic dynamics for multi-agent systems,'' Physica D: Nonlinear
Phenomena, vol.~240, no.~4-5, pp.~432--442, 2 2011.\\[2pt]
[24] L.M.~Miller, Y.~Silverman, et al., ``Ergodic Exploration of
Distributed Information,'' IEEE Transactions on Robotics, vol.~32,
no.~1, pp.~36--52, 2 2016.\\[2pt]
[25] L.M.~Miller and T.D.~Murphey, ``Trajectory optimization for
continuous ergodic exploration,'' in Proceedings of the American
Control Conference, 2013.\\[2pt]
[26] C.~Bilaloglu, T.~L\"ow, and S.~Calinon, ``Whole-body ergodic
exploration with a manipulator using diffusion,'' IEEE Robotics and
Automation Letters, vol.~8, no.~12, pp.~8581--8587, 2023.\\[2pt]
[27] V.L.~Popov and M.~He\ss, Method of dimensionality reduction in
contact mechanics and friction. Springer, 2015.\\[2pt]
[28] L.~Beber, E.~Lamon, et al., ``Force-based viscosity and
elasticity measurements for material biomechanical characterisation
with a collaborative robotic arm,'' IEEE Transactions on
Instrumentation and Measurement, pp.~1--1, 2025.\\[2pt]
[29] P.N.~Wells and H.-D.~Liang, ``Medical ultrasound: imaging of
soft tissue strain and elasticity,'' Journal of the Royal Society
Interface, vol.~8, no.~64, pp.~1521--1549, 2011.
\end{footnotesize}

\end{document}
